\section{A diagonal-gauge obstruction and the full signed source cover}
\label{sec:positive-source-cover}

This section concerns the original honest source basis of
$\check X_{(7,5,3)}$, not an unspecified prospective canonical basis.
It gives a pointwise obstruction for every real $q>0$, including $q=1$,
and an integral positive cover with an exact quotient to all four source
generators. The cover's nonzero quantum-relation defect is also computed.
The ordinary sign-readjustment obstruction is already stated in
GCT~IV~\cite{blasiakmulmuleysohoni2013gct4}, immediately before its
Example \texttt{ex not positive}. The additional statements here are
the arbitrary diagonal-gauge invariant, its precise effectiveness
consequence, and the full source cover and defect maps.

\subsection{Original coefficients and a pointwise invariant}

Retain $A=\mathbb Z[q,q^{-1}]$, $K=\mathbb Q(q)$ and
$[m]=\sum_{j=0}^{m-1}q^{m-1-2j}$ for positive integers $m$.
Put $[0]=0$ and $[-m]=-[m]$. The fixed basis of the honest lattice
$X_A$ is $b_0,\ldots,b_{1259}$. In the original source indexing,
\[
 T=b_{126},\qquad T_1=b_{146},\qquad T_2=b_8,\qquad Q=b_{144}.
\]
Their actual column expressions, with the original scaling, are
\begin{align*}
 T&=(123)(123)(124)(12)(12)(1)(1),\\
 T_1&=(123)(123)(124)(23)(12)(1)(1),\\
 T_2&=-[2]^{-1}(123)(123)(123)(12)(12)(4)(1),\\
 Q&=-[2]^{-1}(123)(123)(124)(13)(12)(4)(1).
\end{align*}
The quotient calculation of the original words supplies these identities
inside $X_A$; the displayed $[2]^{-1}$ is part of the representative,
not an inversion newly made in $A$. The $W$-weight of each vector is
$(9,6)$. Their $V$-weights are, in order, $(12,3),(11,4),(11,4),(10,5)$.
The printed endpoint with a $(23)$ in place of $(13)$ has $W$-weight
$(8,7)$ and is a different vector; it is not used as $Q$.

The following are entries in the full $1260$ by $1260$ matrix $F_V$,
before any four-state quotient:
\[
\begin{array}{c|c|c}
\text{source}&\text{target}&\text{coefficient}\\\hline
126&146&\alpha=[7]-[3]=q^6+q^4+q^{-4}+q^{-6}\\
126&8&\beta=[8]\\
146&144&\gamma=[6]\\
8&144&-d=-[3]
\end{array}
\]
Other full-matrix entries are retained. In particular the argument
below does not require these four vectors to form a full-group
submodule. Define
\[
 p=[2],\qquad C=q^{10}-q^4-q^{-4}+q^{-10}.
\]
Expansion of the finite Laurent sums gives
\begin{equation}\label{eq:positive-source-path-identity}
 \alpha\gamma-\beta d=pC,
 \qquad C=(q-q^{-1})^2[7][3].
\end{equation}
For example $[6]=[3](q^3+q^{-3})$ and direct multiplication gives
$(q^3+q^{-3})([7]-[3])-[8]=[2](q-q^{-1})^2[7]$;
these two identities prove the first formula with no specialization.
The second formula expands to exactly the four terms defining $C$.

Fix a real $t>0$ and evaluate $q=t$. Every term of $[m](t)$ is
positive for $m>0$, and the four displayed terms of $\alpha(t)$ are
positive. Hence the loop invariant of the two original length-two paths is
\begin{equation}\label{eq:positive-source-holonomy}
 \mathcal H(t)=\frac{\alpha(t)\gamma(t)}{-\beta(t)d(t)}
 =-1-\frac{p(t)C(t)}{\beta(t)d(t)}\leq-1.
\end{equation}
Equality holds exactly at $t=1$, since $[7](t)[3](t)>0$ and
$(t-t^{-1})^2=0$ exactly there. Thus the invariant is strictly negative
on the entire positive real axis. At $t=1$ the four coefficients are
$4,8,6,-3$, respectively, and the two products are $24,-24$.
The vanishing of $C(1)$ therefore retains the opposite signs of these
two nonzero paths.

\begin{theorem}\label{thm:positive-source-arbitrary-gauge}
For every $t>0$ there is no diagonal matrix
$D=\operatorname{diag}(d_0,\ldots,d_{1259})$ with
$d_i\in\mathbb C^\times$ such that all entries of
$D^{-1}F_V(t)D$ are nonnegative real numbers. In particular there is
no such simultaneous diagonal gauge for the four original Chevalley
generators. The assertion still holds after a permutation of the source
basis. No diagonal gauge over $\mathbb C(q)$ can make all these entries
nonnegative Laurent polynomials, or rational functions that are
nonnegative wherever defined on a nonempty positive real interval.
\end{theorem}
\begin{proof}
The entry from $i$ to $j$ is multiplied by $d_i/d_j$.
Each of the two path products from $126$ to $144$ is consequently
multiplied by the same nonzero factor $d_{126}/d_{144}$.
Their quotient remains $\mathcal H(t)<0$. If all four transformed
entries were nonnegative real numbers, they would be strictly positive
because none vanishes, so the quotient would be positive. This is a
contradiction even when the $d_i$ are complex. A permutation merely
relocates the same four entries. For a rational diagonal gauge, choose
$t$ in the proposed interval away from the finite set of poles and
zeros of its diagonal entries. The preceding pointwise contradiction
then applies. A nonzero Laurent polynomial with nonnegative
coefficients is strictly positive at every positive real $t$, which
proves the polynomial assertion as a special case.
\end{proof}

\subsection{The exact sign class, and categorical effectiveness}

For a family of real matrices with fixed nonzero support, form an
undirected multigraph: each nonzero matrix entry gives a labelled edge
between its source and target, retaining its matrix colour and sign.
Loops and parallel edges are retained. Let the edge sign be
$\sigma_e\in\{1,-1\}$. A traversal in either direction uses the same
sign. A closed walk has sign equal to the product of these edge signs.
In additive notation $\sigma_e=(-1)^{\epsilon_e}$, the assignment
$\epsilon$ is a $\mathbb Z/2$-valued graph one-cochain. Since a graph
has no two-cells, its class modulo vertex coboundaries is an element
of $H^1(\Gamma;\mathbb Z/2)$, whose definition here is precisely that
quotient of edge assignments.

\begin{lemma}\label{lem:positive-source-switching}
All entries of a family can be made nonnegative by a common diagonal
complex gauge if and only if every closed walk of its sign graph has
positive sign. When this holds, a diagonal gauge with entries in
$\{1,-1\}$ already suffices.
\end{lemma}
\begin{proof}
Write a putative complex diagonal entry as $d_i=r_i z_i$, where
$r_i>0$ and $|z_i|=1$. A real nonzero coefficient of sign $\sigma_e$
from $i$ to $j$ becomes positive only if
$\sigma_e z_i/z_j=1$, or $z_j=\sigma_e z_i$.
Multiplying these identities around a closed walk gives product sign
one. Conversely, in each connected component choose a root and set
its sign to one. Define the sign of a vertex to be the product of edge
signs along any root-to-vertex walk. Two choices have quotient the
sign of a closed walk, so the value is independent of the walk. The
endpoint signs then satisfy $s_j=\sigma_e s_i$ on every edge.
The diagonal entries $s_i$ make each coefficient positive. This also
proves that the obstruction is exactly $[\epsilon]$, not merely a
necessary condition extracted from one failed gauge.
\end{proof}

All $8162$ nonzero entries of the four original matrices are
coefficientwise sign coherent, in agreement with GCT~IV's theorem
\texttt{t positivity}. Thus their signs do not vary with $t>0$.
The exact census in the original basis is
\[
\begin{array}{c|r|r|r}
 &\text{positive entries}&\text{negative entries}&\text{total}\\\hline
 F_V&1648&350&1998\\
 E_V&1837&246&2083\\
 F_W&1848&150&1998\\
 E_W&1963&120&2083
\end{array}
\]
The graph is connected. Its $8162$ labelled edges give cycle rank
$8162-1260+1=6903$; it has $6146$ unordered vertex pairs after
parallel edges are identified. These finite statements are certified
by a spanning traversal of every original matrix entry, retained in
the accompanying independent graph certificate.

The negative cycle displayed in
\eqref{eq:positive-source-holonomy} has minimum possible length.
Indeed every edge changes exactly one of the first weight coordinates
by one. It therefore reverses the parity of their sum, so there are no
loops or odd cycles. Direct comparison of all parallel entries gives
the same sign for every edge on a fixed unordered pair, so a two-edge
closed walk has positive sign. Our four-edge walk is negative, proving
minimality. An additional two-colour witness uses the actual entries
\[
 (F_V)_{20,0}=[7],\quad (F_W)_{20,10}=[5],\quad
 (F_W)_{136,10}=[3],\quad (F_V)_{136,0}=-[3].
\]
The oriented loop $0\to20\leftarrow10\to136\leftarrow0$ has
coefficient quotient $-[7]/[5]<0$ for $t>0$, and value $-7/5$ at
$t=1$. Cancellation in a divided-power endpoint is not needed for
this second witness.

For an exact formulation of effectiveness, call a labelled graded
category effective when its object classes are finite sums
$\sum_{i,w}n_{i,w}q^w[B_i]$ with $n_{i,w}\in\mathbb Z_{\geq0}$,
and these labels and shifts freely generate its split Grothendieck
group. This includes finite direct sums of labelled graded free
groups, and a graded semisimple category with the stated simples.
An additive endofunctor compatible with the internal grading shifts
and taking objects to objects then has an $A$-linear matrix
over $\mathbb Z_{\geq0}[q,q^{-1}]$: its $i$-th column is the
decomposition of the image of $B_i$.

\begin{corollary}\label{cor:positive-source-effective}
There is no such effective category and grading-shift-compatible
additive source-generator
functor whose induced $F_V$ is identified with the original $F_V$
by $[B_i]\mapsto d_i(q)b_i$ for invertible diagonal rational factors.
This remains impossible after arbitrary internal shifts, parity
changes of individual source labels, or permutations of those labels,
provided the images of objects are required to remain effective in
the resulting labels. At a fixed $t>0$ the coordinate simplicial cone
generated by any individually rescaled source basis is not preserved
by all the generator images.
\end{corollary}
\begin{proof}
An effective matrix evaluates to a nonnegative matrix at each positive
real $t$. Apply Theorem~\ref{thm:positive-source-arbitrary-gauge} at
a nonsingular value. Internal shift and individual parity change
give diagonal factors $q^{w_i}$ and $(-1)^{a_i}$, covered by the
same theorem. The cone assertion is the real matrix statement:
a linear map preserves the coordinate nonnegative cone exactly when
every one of its columns lies in that cone.
\end{proof}

This is an effectiveness obstruction with an explicitly stated
basis constraint. It is not a statement that a Frobenius trace must
be positive: an alternating trace retains its cohomological signs.
Nor does it contradict the nonstandard positivity formulations that
explicitly retain signs and powers of $q-q^{-1}$.

\subsection{An integral positive cover of every original generator}

For each of the four full source matrices $M$, split its entries
coefficientwise as $M=M^+-M^-$ with
$M^+,M^-\in\operatorname{Mat}_{1260}(\mathbb Z_{\geq0}[q,q^{-1}])$
and disjoint coefficient supports. Define
\begin{equation}\label{eq:positive-source-cover-matrices}
 \widetilde X_A=X_A^+\oplus X_A^-,\qquad
 \widetilde M=\begin{pmatrix}M^+&M^-\\M^-&M^+\end{pmatrix},
 \qquad |M|=M^++M^-.
\end{equation}
Here $b_i^+,b_i^-$ are new labelled copies of every original vector;
the signs in their names specify the projection below. No change is
made to the underlying original $b_i$. The formula uses all the
original entries and all their exponents. Since they are sign coherent,
an original edge of positive sign stays on its sheet and an edge of
negative sign switches sheets, always with its original absolute
Laurent coefficient.

Define $A$-linear maps
\[
 \iota:X_A\longrightarrow\widetilde X_A,\quad
       x\longmapsto(x,x),\qquad
 \pi:\widetilde X_A\longrightarrow X_A,\quad
       (x,y)\longmapsto x-y.
\]
The following exact diagram holds for each of the four generators:
\begin{equation}\label{eq:positive-source-cover-exact}
 0\longrightarrow (A^{1260},|M|)
 \xrightarrow{\ \iota\ }(A^{2520},\widetilde M)
 \xrightarrow{\ \pi\ }(X_A,M)\longrightarrow0.
\end{equation}
To prove this, $\iota$ is injective, $\pi(x,0)=x$, and
$\pi(x,y)=0$ exactly when $x=y$. Direct block multiplication gives
\[
 \pi\widetilde M=M\pi,\qquad
 \widetilde M\iota=\iota|M|.
\]
The underlying sequence of free $A$-modules is split by
$s(x)=(x,0)$, with no division by two; the section is not asserted
to commute with the generators. Thus there is neither an unmentioned
$2$-inversion nor an integral torsion kernel.

Give both copies of $b_i$ its original four weight coordinates.
Every original Cartan diagonal matrix has Laurent monomial entries,
so it lifts to the same diagonal matrix on both sheets. Its inverse
lifts in the same way. The two original total degree-fifteen central
characters remain $q^{15}$. The displayed inclusion and projection
commute with these diagonal operators. The Cartan conjugation rules
hold on the cover because each nonzero lifted edge has exactly the
original source and target weights. In particular, the sign cover
does not replace a raising edge by a lowering edge.

To describe its exact algebraic scope, let $\mathcal T$ be the free
associative $A$-algebra on the four Chevalley symbols, the original
Cartan symbols and their named inverses, and the integral Cartan
operators $H_V,H_W$ defined below. On the two sheets and on the
kernel these last operators act by the same original diagonal
matrices. Acting by the specified
matrices makes all three terms of
\eqref{eq:positive-source-cover-exact} actual $\mathcal T$-modules,
and the displayed maps are $\mathcal T$-linear. For any polynomial
relation $r\in\mathcal T$ whose original source matrix is zero,
\[
 \pi r(\widetilde M)=0.
\]
Consequently there is a unique, explicitly defined $A$-linear defect
map $\Delta_r:\widetilde X_A\to A^{1260}$ with
\begin{equation}\label{eq:positive-source-relation-defect}
 r(\widetilde M)=\iota\Delta_r,
 \qquad \Delta_r=\operatorname{pr}_+r(\widetilde M).
\end{equation}
Existence follows from the equality of the plus and minus output
coordinates, and uniqueness from injectivity of $\iota$. If
$r(\widetilde M)\iota$ is evaluated instead, the same identities give
$\Delta_r\iota=r(M_{\ker})$, where $M_{\ker}$ denotes the four
absolute Chevalley matrices and the original Cartan diagonals.
Thus the complete relation loss is an actual
map into the displayed kernel. The quotient by that kernel recovers
all original source relations.

This doubled graph is the covering determined by the sign class.
It is connected: paths from a root lift to one sheet at any target;
a negative closed walk switches the root sheet, and then the same
paths reach the other sheet at every target. A vertex sign on the
cover, positive on the plus sheet and negative on the minus sheet,
trivializes the pullback of the sign class. Conversely a one-sheet
cover cannot trivialize its nonzero class. A finite graph covering
of a connected graph has a constant number of vertices in each
fibre: lift any path and its inverse to obtain bijections of fibres.
Therefore this degree-two, $2520$-vertex cover has the least possible
degree among nonempty graph covers that trivialize this sign class. This
minimality statement is about such graph covers, not arbitrary
changes of representation basis or unrestricted categorical models.

\subsection{A nonzero full-source quantum-relation defect}

Let $H_V$ be the original diagonal operator with entry
$[w_{V,1}(b_i)-w_{V,2}(b_i)]$ on $b_i$; equivalently over $K$ it is
$(K_V-K_V^{-1})/(q-q^{-1})$. Its integral diagonal formula, with its
signs, defines the operator already over $A$. Let
$\widetilde H_V=\operatorname{diag}(H_V,H_V)$, the actual same
Cartan expression on both sheets.

The exact source matrices give
\begin{equation}\label{eq:positive-source-nonzero-defect}
 (\widetilde E_V\widetilde F_V-
   \widetilde F_V\widetilde E_V-\widetilde H_V)b_{132}^+
       =[2][3](b_{148}^++b_{148}^-).
\end{equation}
Here the original source labels are
\[
 b_{132}=(123)(123)(124)(12)(12)(3)(3),\qquad
 b_{148}=(123)(123)(124)(12)(23)(1)(3).
\]
Both have $V$-weight $(10,5)$ and $W$-weight $(9,6)$.
For completeness, all paths from $132$ to $148$ contributing to
the two compositions are the following; coefficients are recorded
in traversal order:
\[
\begin{array}{c|r|c|c}
\text{composition}&\text{intermediate index}&\text{first coefficient}
                                         &\text{second coefficient}\\\hline
 E_VF_V&152&-[3]&[2]\\
 E_VF_V&158&[7]&[2]\\
 F_VE_V&128&[2]&[7]-[3]
\end{array}
\]
On the original source their sum is
$-[3][2]+[7][2]-[2]([7]-[3])=0$.
On the cover the first path ends on the minus sheet and the second
on the plus sheet. The subtracted third path ends on the plus sheet.
Each output sheet therefore has coefficient $[2][3]$. The full
column calculation verifies that all other target entries in the
commutator defect vanish, proving
\eqref{eq:positive-source-nonzero-defect}. The diagonal term cannot
affect target $148\ne132$. The original full matrices, including
the same diagonal, satisfy the zero commutator relation in every
column; thus this is precisely a failure upstairs, with zero image
under $\pi$. Its coefficient expands to
\[
 [2][3]=q^3+2q+2q^{-1}+q^{-3},
\]
which is nonzero over $A$ and strictly positive at every $t>0$.
The $W$-coloured calculation gives
\[
 (\widetilde E_W\widetilde F_W-
   \widetilde F_W\widetilde E_W-\widetilde H_W)b_{256}^+
       =[2][3](b_{263}^++b_{263}^-).
\]
Each colour has exactly $52$ even-sheet input columns with nonzero
commutator defect. These claims use all target coordinates, not a
principal submatrix. The independent Laurent checker verifies them
column by column. It follows that the cover in
\eqref{eq:positive-source-cover-matrices} is not a full quantum-group
representation, although its exact quotient is the original one.

The failure already includes the commuting lowering generators:
\begin{equation}\label{eq:positive-source-mixed-defect}
 [\widetilde F_V,\widetilde F_W]b_0^+
       =-[3]^2(b_{146}^++b_{146}^-),
 \qquad b_0=(123)(123)(123)(12)(12)(1)(1).
\end{equation}
All contributing original paths from $0$ to $146$ are
\[
\begin{array}{c|r|c|c}
\text{composition}&\text{intermediate index}&\text{first coefficient}
                                         &\text{second coefficient}\\\hline
 F_VF_W&126&[3]&[7]-[3]\\
 F_WF_V&20&[7]&[3]\\
 F_WF_V&136&-[3]&[3]
\end{array}
\]
The original difference is $[3]([7]-[3])-[7][3]+[3]^2=0$.
On the cover the third path ends on the minus sheet and is
subtracted with its positive absolute coefficient. The plus-sheet
difference is also $-[3]^2$. Every other target entry cancels,
as the full original-path enumeration verifies. In total the
numbers of nonzero even-sheet input columns for the six
commutator defects are
\[
\begin{array}{c|rrrrrr}
\text{commutator}&[E_V,F_V]&[E_W,F_W]&[E_V,E_W]&[E_V,F_W]
                  &[F_V,E_W]&[F_V,F_W]\\\hline
\text{columns}&52&52&68&52&52&50
\end{array}
\]
For the first two columns the stated Cartan diagonal is subtracted;
for the four mixed columns the original commutator is zero.
The independent checker verifies all six original identities and
all their exact diagonal-kernel defects over $A$.

\subsection{An actual categorical map and the retained cone}

Let $\mathcal G$ be the category of finite free internally graded
abelian groups with finite support. A Laurent polynomial
$a(q)=\sum_w a_wq^w$ with nonnegative integer coefficients defines
the actual group $L_a=\bigoplus_w\mathbb Z^{a_w}\langle w\rangle$.
Use the ordered basis labels $(w,1),\ldots,(w,a_w)$ and set $L_0=0$.
Define an additive functor on $\mathcal G^{2520}$ by
\[
 (\widetilde{\mathsf M}V)_i
       =\bigoplus_j L_{\widetilde M_{ij}}\otimes_{\mathbb Z}V_j.
\]
On morphisms use the identity on each first tensor factor. Its
Grothendieck matrix is exactly $\widetilde M$, since tensor products
add degrees and multiply ranks. This formula realizes every one of
the positive-cover coefficients as actual objects, over $\mathbb Z$.

Let $\mathcal S$ be internally graded finite free super abelian
groups, with even morphisms, and let $\Pi$ exchange parity.
It satisfies $\Pi^2=1$ on objects and maps. Its super Grothendieck
group is defined by direct-sum relations and $[\Pi V]=-[V]$;
the signed rank character identifies that group with $A$.
The parity-summing functor
\[
 \mathsf Q:\mathcal G^{2520}\to\mathcal S^{1260},\qquad
 (V_i^+,V_i^-)_i\longmapsto(V_i^+\oplus\Pi V_i^-)_i
\]
has induced map $\pi$, with every plus component placed in even
parity before taking the direct sum. Define
\[
 (\mathsf M W)_i=\bigoplus_j
   \left(L_{M^+_{ij}}\otimes W_j\right)
   \oplus\left(L_{M^-_{ij}}\otimes\Pi W_j\right).
\]
There is a natural parity-preserving isomorphism
\begin{equation}\label{eq:positive-source-super-map}
 \mathsf Q\widetilde{\mathsf M}\simeq\mathsf M\mathsf Q.
\end{equation}
To verify it, expand the two sides. The even output summands on
each side are $L_{M^+_{ij}}\otimes V_j^+$ and
$L_{M^-_{ij}}\otimes V_j^-$; the odd output summands are
$L_{M^+_{ij}}\otimes V_j^-$ and
$L_{M^-_{ij}}\otimes V_j^+$. Match identical labelled tensor
factors by the identity. This specifies every component, its inverse,
and naturality. Thus the signed source action has an actual integral
super-object correspondence with the full effective cover, not only
an equality of matrix dimensions.

The quantum relations have not been imposed as functor isomorphisms
by this definition. The nonzero defect just computed precludes that
interpretation for the positive cover. The full generators on
super Grothendieck groups do satisfy the source relations, through
\eqref{eq:positive-source-super-map} and the proven matrix action.
No assertion of a canonical-basis convolution model follows merely
from the existence of these additive functors.

There is also a precise relation to the previously constructed
four-state Frobenius cone. Let $\mathsf A,\mathsf B,\mathsf G,\mathsf J$
be the graded groups with respective exponent sets
\[
\begin{split}
 &\{-6,-4,4,6\},\qquad\{-7,-5,-3,-1,1,3,5,7\},\\
 &\{-5,-3,-1,1,3,5\},\qquad\{-2,0,2\}.
\end{split}
\]
The two lifted paths from $T^+$ end in $Q^+$ with
$\mathsf P=\mathsf G\otimes\mathsf A$ and in $Q^-$ with
$\mathsf R=\mathsf J\otimes\mathsf B$.
Applying $\mathsf Q$ gives $\mathsf P\oplus\Pi\mathsf R$,
with character $\alpha\gamma-d\beta=pC$ and with the intermediate
labels $T_1,T_2$ still attached. The actual integral map
$\theta:\mathsf R\to\mathsf P$ used in the Frobenius cone matches,
in each weight, equal-position basis pairs in lexicographic order.
Its twenty unit entries cancel twenty matched pairs; the cone has
four residual even and four residual odd basis vectors. This
cross-path differential is additional map data, not the differential
of the split path object. The cover and the parity map therefore
recover exactly its two branch objects before that differential is
installed. More explicitly, place an arbitrary input $V$ at $T^+$
and project the square functor to the two labelled output sheets
at $Q$. The extracted endpoint functors are
$V\mapsto\mathsf P\otimes V$ and
$V\mapsto\mathsf R\otimes V$. The map
$\theta\otimes1_V:\mathsf R\otimes V\to\mathsf P\otimes V$
is natural, and its cone is exactly the specified endpoint cone
functor. Its underlying graded parity object is the image under
$\mathsf Q$ of the two extracted paths. The identity of those
graded objects is not claimed to be a chain map from the split
zero-differential path complex to the cone.

\subsection{What changes when the basis is allowed to mix}

For each fixed $t>0$ an effective cone for the four full generators
does exist after an explicit type of nondiagonal source change.
Here the exact source commutator is read as the integral formula
$[E,F]|_{X_h}=[h](t)$, including its value at $t=1$.
This assertion can be proved without assuming a positive basis in
the original labels.

Consider one colour on any of the present finite-dimensional
specializations. If $z$ is highest of weight $n$, the commutator
gives by induction
\[
 E F^jz=[j](t)[n-j+1](t)F^{j-1}z.
\]
The induction uses the finite-sum identity
$[j+1][n-j]=[j][n-j+1]+[n-2j]$, verified by multiplying the
Laurent sums, so it also holds at $t=1$.
Let $k$ be the largest integer with $F^kz\ne0$. Applying the
formula to $F^{k+1}z=0$ gives
$[k+1](t)[n-k](t)F^kz=0$. For integral $h$, $[h](t)=0$ exactly
when $h=0$. Thus $k=n\geq0$. It follows that the divided string
$z_i=F^iz/[i]!(t)$, $0\leq i\leq n$, has all its stated vectors
nonzero and satisfies
\[
 Fz_i=[i+1](t)z_{i+1},\qquad
 Ez_i=[n-i+1](t)z_{i-1},
\]
where the out-of-range terms vanish. All nonzero coefficients are
strictly positive.

To prove these strings span the entire representation, take its
largest occurring weight $n$ and a basis of its highest space.
Their strings are independent: different levels have different
weights, and applying the corresponding power of $E$ to a relation
at one level reduces it to a relation in that highest basis with a
nonzero common coefficient. Their sum $N$ is a submodule. Induct
on dimension in the quotient, whose largest weight is smaller.
For a quotient highest vector of weight $m$, lift it to $z$ of
weight $m$. Then $Ez\in N_{m+2}$. On every top string,
$E:N_m\to N_{m+2}$ is surjective when the target occurs, because
the coefficient just displayed is nonzero; the quotient string
has $m\geq0$, so there is no endpoint exception. Subtract a
preimage to make the lift highest. Its length-$m$ string maps
isomorphically to the selected quotient string. These lifted
strings are independent modulo $N$ and split off the quotient.
This completes the induction.

Apply the construction first to the $V$ generators. The commuting
$W$ generators preserve each $V$ highest subspace, where the same
argument applies. If $H_{n,m}(t)$ is the simultaneous highest
subspace of weight ratios $(n,m)$ and a real basis of it is chosen,
define $L_n(t)$ to have the divided-string basis $e_i$, $0\leq i\leq n$,
the displayed $E,F$ action, and original $\mathfrak{gl}_2$ weights
\[
 \left(\frac{15+n}{2}-i,\frac{15-n}{2}+i\right).
\]
The occurring $n$ are odd, since all original total weights are
fifteen, so these coordinates are integers. Define $L_m(t)$ in
the same way for $W$. Then the map
\[
 \Phi_t:\bigoplus_{n,m} L_n(t)\otimes L_m(t)\otimes H_{n,m}(t)
       \longrightarrow X_A\otimes_{q\mapsto t}\mathbb R,
 \quad e_i\otimes f_j\otimes z\longmapsto
              F_V^{(i)}F_W^{(j)}z
\]
is therefore an isomorphism. The multiplicity factor has trivial
quantum-group action in the domain. The formulas above prove that
the map intertwines all raising and lowering generators; the
original weight coordinates prove the Cartan and degree-fifteen
central characters as well. This is the specialization of the
actual string construction, now justified at every positive $t$.

The nonnegative span of these string vectors is a simplicial cone
preserved by all four generators and their Cartan diagonals: their
matrices have nonnegative entries in these strings, and each Cartan
eigenvalue is a positive power of $t$. Its matrix $\Phi_t$ is
obtained from actual simultaneous highest kernels and the original
divided powers, so it is an exact correspondence with the source.
Theorem~\ref{thm:positive-source-arbitrary-gauge} proves that
$\Phi_t$ cannot be a diagonal matrix times a permutation. Thus the
fixed-source obstruction and this effective mixed-basis cone are
related by an explicit proved intertwiner. The construction depends
on choices in the highest multiplicity spaces and is not asserted
to be the source-compatible canonical geometric realization sought
in the nonstandard Riemann-hypothesis programme.

\paragraph{Reproduction and scope.}
The complete original matrices and basis, with their original
coefficient and word data, are the frozen inputs in
\texttt{agents/full\_source\_generators}. The accompanying
\texttt{check\_positive\_cover.py} verifies their hashes, every edge
sign, all four retained coefficients, both Laurent identities, every
column of the inclusion and projection intertwiners, the two full
original commutator identities, and every cover commutator defect.
The independent sign audit checks graph connectivity, the cycle
class, the shortest-cycle assertion and the same retained entries.
The primary-source loci are GCT~IV TeX lines $7240$--$7323$
for sign coherence, $7324$--$7326$ for the already stated
sign-readjustment obstruction, and $7328$--$7350$ for the printed
example and its two paths; its frozen source hash is recorded in
the machine receipt. The new obstruction is to a stated effective
realization preserving the original basis up to diagonal gauge.
The exact cover and mixed-string maps do not prove a nonstandard
Riemann hypothesis, an asymptotic class-variety obstruction, or a
Boolean or algebraic complexity-class separation.
